\subsection{9.1 Was der vorhandene Vierzustandszeuge bereits
leistet}\label{was-der-vorhandene-vierzustandszeuge-bereits-leistet}

Im wirklichen Tensor \(W\) existieren die Paarungen \((4,57)\) in Kanal
0 und \((4,58)\) in Kanal 1. Mit den sieben Pauli-Blockern

\[S_b=\{8,16,28,32,44,52,56\}\]

und einem zusätzlich gewährten Bosonmischer schließt die vollständige
Paarwechselwirkung auf den vier N=9-Zuständen

\[
|1\rangle=S_b+\{4,57\},\quad |2\rangle=S_b+b_0,
\quad |3\rangle=S_b+b_1,\quad |4\rangle=S_b+\{4,58\}.
\]

Nach Phasenwahl ist der exakte Block

\[H_4=\begin{pmatrix}
0&g&0&0\\g&\Delta&J&0\\0&J&\Delta&g\\0&0&g&0
\end{pmatrix}.\]

Die Quellen beweisen seine Invarianz unter sämtlichen nativen
Paartermen; es handelt sich nicht um eine willkürliche
Vierzustandsabschneidung. Für

\[g/\Delta=1/20,\quad J/\Delta=1-1/(10\sqrt3),
\quad t\Delta=20\pi\sqrt3\]

gilt

\[P_{1\to4}>\frac{2009992727}{2022609600}>0.99376208.\]

Der Mischer verändert allerdings eine erhaltene innere Cartanladung mit
Werten \((7,7,9,9)\). Er ist aus dem bisherigen diese Ladung erhaltenden
Alphabet nicht erzeugbar. Die native Teilchenladung \(N=9\) ist eine
andere Größe und bleibt schon im ursprünglichen Block erhalten. {[}S033,
§5{]}

\subsection{9.2 Explizite ladungserhaltende
Erweiterung}\label{explizite-ladungserhaltende-erweiterung}

\textbf{Neue Konstruktion, bedingt.} Ergänze eine zweidimensionale
Referenz mit Zuständen \(|+\rangle,|-\rangle\) und

\[Q_R|\pm\rangle=\pm|\pm\rangle,\qquad
R_-=|-\rangle\langle+|.\]

Schreibe \(Q_S=\operatorname{diag}(7,7,9,9)\) und
\(M=|3\rangle\langle2|\). Dann gilt

\[[Q_S,M]=2M,\qquad [Q_R,R_-]=-2R_-.\]

Für \(H_0=H_4|_{J=0}\) definiere

\[
\boxed{H_{\mathrm{rel}}=H_0\otimes I
+J(M\otimes R_-+M^\dagger\otimes R_-^\dagger)
+\kappa I\otimes Q_R^2.}
\]

Dieser achtdimensionale Hamiltonoperator ist selbstadjungiert und
erfüllt exakt

\[[H_{\mathrm{rel}},Q_S\otimes I+I\otimes Q_R]=0.\]

Der Raum

\[\mathcal K=\operatorname{span}\{
|1,+\rangle,|2,+\rangle,|3,-\rangle,|4,-\rangle\}\]

hat durchgehend Gesamtladung acht. Er ist invariant, und in dieser Basis
gilt

\[H_{\mathrm{rel}}|_{\mathcal K}=H_4+\kappa I_4.\]

\textbf{Beweis.} Die beiden nativen Endkopplungen ändern \(Q_S\) nicht.
Der mittlere Übergang erhöht \(Q_S\) um zwei und erniedrigt gleichzeitig
\(Q_R\) um zwei. Keine der drei Kopplungen verlässt die vier angegebenen
Zustände. Wegen \(Q_R^2=I\) ist der Referenzenergieterm dort skalar. Die
Einschränkung ist daher genau die behauptete Matrix. Alle Identitäten
wurden symbolisch geprüft. \(\square\)

Damit wird die vorhandene Transferwahrscheinlichkeit unverändert auf die
erweiterte Konstruktion übertragen. Der Endzustand ist nun
\(|4,-\rangle\) statt \(|4,+\rangle\). Die mitgerechnete Referenz
erklärt exakt, wohin die fehlende Cartanladung geht.

\begin{figure}
\centering
\pandocbounded{\includegraphics[keepaspectratio,alt={Bedingter relationaler Transfer. Links: numerische Zeitkurve des exakt definierten Vierzustandsblocks; die rigorose Endpunktschranke stammt aus S033. Rechts: die neue exakte Ladungsbilanz. Die Referenz und ihre gekoppelte Operation sind zusätzliche Ressourcen.}]{figures/02_relationaler_transfer.pdf}}
\caption{Bedingter relationaler Transfer. Links: numerische Zeitkurve
des exakt definierten Vierzustandsblocks; die rigorose Endpunktschranke
stammt aus S033. Rechts: die neue exakte Ladungsbilanz. Die Referenz und
ihre gekoppelte Operation sind zusätzliche Ressourcen.}
\end{figure}

\subsection{9.3 Was damit fundamental gelöst wird -- und was
nicht}\label{was-damit-fundamental-geluxf6st-wird-und-was-nicht}

Gelöst ist die konkrete mathematische Frage, ob die festgestellte
Cartan-Schranke einen Transfer auch bei mitgeführtem Ladungsausgleich
grundsätzlich verbietet. Das tut sie nicht: Die Erweiterung liefert
einen expliziten positiven Zeugen. Sie bewahrt die betrachtete
Gesamt-Cartanladung, ohne den ursprünglichen isolierten Mischer als
bereits verfügbar auszugeben.

Die Konstruktion leitet weder die Referenzzustände noch die gekoppelte
Operation aus TFPT ab. Sie behauptet auch keine Erhaltung der
vollständigen inneren Gruppe durch diese Erweiterung. Dafür müssten
vollständige Darstellungen, alle anderen Ladungen und passende
Intertwiner ergänzt und geprüft werden. Die Referenz wurde lediglich für
die einzelne konkret blockierende Cartanladung gebaut. Ihr
physikalischer Träger und ihre sonstigen Quantenzahlen sind offene
Eingaben.

Allgemein kann eine Operation \(M_q\) mit \([Q_S,M_q]=qM_q\) durch eine
Referenzoperation \(R_{-q}\) mit \([Q_R,R_{-q}]=-qR_{-q}\) zu einem
neutralen Produkt ergänzt werden. Dieses Prinzip ist in der Theorie von
Quantenreferenzen und Superselektion bekannt. Hier wird es auf den
vorhandenen nativen Transferblock angewandt und vollständig
ausgerechnet. {[}E4{]}

\subsection{9.4 Die Referenz wird verbraucht: ein einfacher
Unmöglichkeitssatz}\label{die-referenz-wird-verbraucht-ein-einfacher-unmuxf6glichkeitssatz}

Beginnt die Referenz in einer scharfen Ladung \(r\) und erhöht sich die
Systemladung in einem erfolgreichen Zweig um \(q\ne0\), muss die
Referenz oder eine weitere Umgebung im selben Zweig ihre Ladung um
\(-q\) ändern. Ein vollständiger Rücklauf zur exakt gleichen scharfen
Referenzladung bei sonst unverändertem Außenraum ist wegen
Gesamtladungserhaltung unmöglich.

Für die obige Konstruktion sind \(|+\rangle\) und \(|-\rangle\)
orthogonal. Eine zweistufige Referenz ist für einen solchen scharfen
einmaligen Ausgleich minimal; ein eindimensionaler Träger kann keinen
nichtverschwindenden Ladungswechsel aufnehmen. Eine unveränderte
Wiederverwendung verlangt Rücktransport, ein weiteres Reservoir oder
einen anderen Referenzvertrag. Dies ist keine kostenlose katalytische
Operation.

Damit werden Zustandspräparation, Aufzeichnung und Transport miteinander
verbunden: Der Ausgang enthält eine reale Änderung einer mitgerechneten
Ressource. Bei Superpositionen kann sie zudem welche-Zweig-Information
tragen. Wer sie wegtraziert, muss den Verlust von Systemkohärenz in die
Prozessrechnung aufnehmen.

\section{10. Eigene Verbindung III: Überlappung, Dynamik und
geometrische
Krümmung}\label{eigene-verbindung-iii-uxfcberlappung-dynamik-und-geometrische-kruxfcmmung}

\subsection{10.1 Eine gemeinsame Quelle ist nicht automatisch
Bewegung}\label{eine-gemeinsame-quelle-ist-nicht-automatisch-bewegung}

Sei \(T:\mathbb C^m\to\mathcal H\) ein Adapter, dessen Spalten die
betrachteten Zustände sind. Dann

\[S=T^\dagger T,\qquad K=T^\dagger HT.\]

Ist \(S>0\), erhält man mit \(J=TS^{-1/2}\) einen isometrischen Adapter
und

\[h=J^\dagger HJ=S^{-1/2}KS^{-1/2}.\]

Bei singulärem \(S\) wird zuerst sein Nullraum entfernt. Ist
\(K=\epsilon S\), dann ist \(h=\epsilon I\). Offdiagonale Einträge in
\(K\) können also vollständig durch die nichtorthogonale Beschriftung
entstehen. Für tatsächliche Dynamik muss zusätzlich die Kopplung zum
ausgelassenen Raum kontrolliert werden; Abschnitt 8.5 liefert dafür
einen einfachen Vertrag.

Auf dem einzelnen entarteten nativen Lochpol gilt genau
\(P_hHP_h=E_hP_h\). Neue Charts innerhalb dieses Polraums ergeben
zwangsläufig \(K=E_hS\). Sie erzeugen durch bloße Überlappung keine
räumliche Ausbreitung. Die neue gemeinsame Quelle muss wirklich mehr als
eine neue Benennung dieser Zustände liefern. {[}S033, §7{]}

\subsection{10.2 Zeitabhängige Rahmen brauchen einen
Verbindungsterm}\label{zeitabhuxe4ngige-rahmen-brauchen-einen-verbindungsterm}

Für einen glatten isometrischen Rahmen \(J(t)\) und \(\psi(t)=J(t)c(t)\)
ergibt die projizierte Schrödingergleichung

\[\boxed{i\dot c=(J^\dagger HJ-iJ^\dagger\dot J)c.}\]

Der zweite Term ist nötig, weil die Koordinaten selbst bewegt werden.
Eine Rechnung, die nur \(J^\dagger HJ\) betrachtet, kann eine
Rahmendrehung fälschlich als physische Dynamik lesen. Bei einem echten
bewegten Unterraum muss die außerhalb liegende Gleichung ebenfalls
erfüllt oder durch eine adiabatische beziehungsweise andere
Fehlerabschätzung kontrolliert werden.

Wählt man \(J(t)=J_0V(t)\) innerhalb eines festen Unterraums, ist die
zusätzliche Struktur reine Rahmenwahl. Eine Vorhersage für einen
festgelegten Detektor darf davon nicht abhängen. Deshalb müssen
Zustände, Generator und Messinstrument zusammen transformiert werden.

\subsection{10.3 Feste vollständige Rahmen und bewegte Projektoren
unterscheiden
sich}\label{feste-vollstuxe4ndige-rahmen-und-bewegte-projektoren-unterscheiden-sich}

Sind \(U_{xy}=V_x^\dagger V_y\) bloße Basiswechsel vollständiger
unitärer Rahmen, teleskopiert eine geschlossene Schleife:

\[U_{xy}U_{yz}U_{zx}=I.\]

Anders ist es für isometrische Einbettungen \(J_x\) echter Unterräume
mit \(P_x=J_xJ_x^\dagger\). Dann

\[J_x^\dagger J_yJ_y^\dagger J_zJ_z^\dagger J_x
=J_x^\dagger P_yP_zJ_x,\]

und dies muss nicht die Identität sein. Die Zwischenprojektionen sind
zusätzliche geometrische beziehungsweise physische Struktur. Die
Überlappungsmatrizen sind im Allgemeinen keine unitären
Transportoperatoren; ihre Beträge und mögliche
Erfolgswahrscheinlichkeiten müssen erhalten bleiben.

Mit der hermiteschen Verbindung \(A=iJ^\dagger dJ\) ist die Krümmung

\[\boxed{F=dA-iA\wedge A
=i\,dJ^\dagger(I-P)\wedge dJ.}\]

\textbf{Herleitung.} Aus \(J^\dagger J=I\) folgt
\(dJ^\dagger J=-J^\dagger dJ\). Einsetzen in \(dA-iA\wedge A\) trennt
den vollen Ausdruck in den Anteil innerhalb und außerhalb von \(P\); der
innere Anteil hebt sich weg. Ist \(P\) konstant und ändern sich nur
seine Rahmen, verschwindet \(F\). Bewegt sich der physische Unterraum,
kann \(F\ne0\) sein. \(\square\)

Das ist die etablierte Berry-/Wilczek--Zee-Struktur, hier als konkrete
Antwort auf das Überlappungsproblem eingesetzt. Sie ist noch kein
dynamisches Yang--Mills-Feld der Raumzeit. {[}E5{]}

\subsection{10.4 Ein vollständig gerechneter kleiner
Zeuge}\label{ein-vollstuxe4ndig-gerechneter-kleiner-zeuge}

Betrachte die normierte Familie

\[u(\theta,\phi)=
\begin{pmatrix}\cos(\theta/2)\\e^{i\phi}\sin(\theta/2)\end{pmatrix},
\qquad P=uu^\dagger.\]

Direkte Rechnung ergibt

\[A_\theta=0,\quad A_\phi=-\sin^2(\theta/2),
\quad F_{\theta\phi}=-\tfrac12\sin\theta.\]

Aus demselben Projektor folgt die reelle Quantenmetrik

\[ds^2=\operatorname{Re}\langle du,(I-P)du\rangle
=\tfrac14(d\theta^2+\sin^2\theta\,d\phi^2).\]

\textbf{Der verbindende Punkt:} Unterscheidbarkeit und geometrische
Phase stammen hier aus demselben komplexen Tangentialprodukt. Die Metrik
ist sein reeller Teil, die Krümmung hängt an seinem antisymmetrischen
imaginären Teil. Es müssen also nicht zwei unabhängige Strukturen frei
eingesetzt werden. Die Parametrisierung des Projektors selbst bleibt
aber eine Eingabe.

\begin{figure}
\centering
\pandocbounded{\includegraphics[keepaspectratio,alt={Derselbe bewegte Rang-eins-Projektor liefert geometrische Krümmung und Zustandsmetrik. Die Achsen sind Modellparameter; sie sind keine bereits abgeleiteten Raumzeitkoordinaten.}]{figures/03_projektorgeometrie.pdf}}
\caption{Derselbe bewegte Rang-eins-Projektor liefert geometrische
Krümmung und Zustandsmetrik. Die Achsen sind Modellparameter; sie sind
keine bereits abgeleiteten Raumzeitkoordinaten.}
\end{figure}

Eine endliche Kontrolle verwendet

\[u_0=(1,0)^T,\quad u_1=(1,1)^T/\sqrt2,
\quad u_2=(1,i)^T/\sqrt2.\]

Das rahmeninvariante Dreiecksprodukt ist

\[\langle u_0,u_1\rangle\langle u_1,u_2\rangle
\langle u_2,u_0\rangle=(1+i)/4.\]

Seine Phase ist \(\pi/4\), sein Normquadrat \(1/8\). Die selektive
Projektionsfolge \(P_1\), dann \(P_2\), dann \(P_0\) auf \(u_0\) hat
entsprechend Wahrscheinlichkeit \(1/8\) und Amplitude \((1-i)/4\). Die
umgekehrte Orientierung erklärt das Vorzeichen der Phase. Ein Normieren
jedes Schritts ohne seine Ausbeute würde diese reale Ressourcenzahl
unterschlagen. Alle angegebenen Identitäten wurden symbolisch geprüft.

\subsection{10.5 Ein bedingter dynamischer
Parent}\label{ein-bedingter-dynamischer-parent}

Der zusätzliche Hamiltonoperator

\[H(\theta,\phi)=\Delta(I-P(\theta,\phi))\]

hat bei festem Parameter eine Lücke \(\Delta\) und einen
Rang-eins-Grundraum \(P\). Bei kontrollierter langsamer
Parameteränderung liefert er die bekannte geometrische Phase, mit
separat zu kontrollierendem adiabatischem Fehler. Er zeigt konstruktiv,
dass ein bewegter ausgewählter Unterraum die falsche „alles
teleskopiert``-Abkürzung umgehen kann.

Er leitet weder die Parameter noch ihre Raumzeitinterpretation aus TFPT
ab. Ein elementarer geometrischer Phasenraum ist auch kein masseloser
Eichboson. Dafür müssten \(P_x\) oder entsprechende Verbindungsdaten
selbst dynamisch werden, aus einer Quelle stammen und eine positive
lokale Wirkung samt geeignetem Grenzverhalten besitzen.

Das ist ein aussichtsreicher \textbf{bedingter Verbindungsansatz}:
Anstelle willkürlich hinzugefügter Eichpotentiale könnte man die
Geometrie quellseitig ausgewählter niedriger Unterräume untersuchen. Der
unmittelbar entscheidende Test lautet, ob der native Operationssatz
einen physisch veränderlichen Projektor erzeugt. Bewegt er nur den
Rahmen eines festen \(P\), ist dieser Kandidat beendet.

\section{11. Wo Raum und echte Wirkung entstehen
könnten}\label{wo-raum-und-echte-wirkung-entstehen-kuxf6nnten}

\subsection{11.1 Innere Darstellung und
Multiplizität}\label{innere-darstellung-und-multiplizituxe4t}

Unter einer kompakten inneren Gruppe zerfällt ein endlicher Hilbertraum
als

\[\mathcal H=\bigoplus_\lambda R_\lambda\otimes\mathcal M_\lambda.\]

Ein gruppeninvarianter Hamiltonoperator hat die Form

\[H=\bigoplus_\lambda I_{R_\lambda}\otimes h_\lambda.\]

Der Raum \(R_\lambda\) trägt die Art einer Anregung;
\(\mathcal M_\lambda\) unterscheidet unabhängige Vorkommen derselben
Art. Bei Multiplizität eins bleibt nur eine Energie. Bei größerer
Multiplizität ist nichttriviale symmetrieerhaltende Dynamik möglich.
Aber ein Multiplizitätsindex ist erst dann ein Ort, wenn lokale
Zugriffe, Komposition und Ausbreitungsstruktur ihn entsprechend
auszeichnen. {[}S003, §4{]}

Für den invariant gewählten nativen Grundzustand und die irreduzible
64er-Entnahmefamilie folgt nach Schurs Lemma sogar für die volle Antwort

\[\langle f_r\Omega,e^{-it(H-E_0)}f_s\Omega\rangle
=\delta_{rs}c(t).\]

Alle Nebenlinien sind in \(c(t)\) enthalten. Genauere Berechnung
derselben Funktion erzeugt keinen räumlichen Richtungsindex. Ein
globaler Träger muss weitere unabhängig adressierbare Vorkommen oder
eine explizit andere Symmetrie- und Zustandsstruktur liefern.

\subsection{11.2 Zwei verschiedene
Kommutanten}\label{zwei-verschiedene-kommutanten}

Der Kommutant der inneren Gruppe beschreibt die mit ihr verträgliche
Dynamik. Der Kommutant des verfügbaren Operationsalphabets beschreibt
dagegen Unterschiede, die dieses Alphabet nicht auflösen kann. Diese
Bezugsalgebren sind verschieden.

Im N=3-Sektor berichten die aktuellen exakten Rechnungen beispielsweise
folgende Kommutantdimensionen:

{\def\LTcaptype{none} % do not increment counter
\begin{longtable}[]{@{}lr@{}}
\toprule\noalign{}
Gewährtes Operationsalphabet & komplexe Kommutantdimension \\
\midrule\noalign{}
\endhead
\bottomrule\noalign{}
\endlastfoot
\(X,N_b\) & 1444233216 \\
zusätzlich nativer Clock & 240742144 \\
stattdessen zusätzlich \(SU(4)\)-Generatoren & 2247168 \\
stattdessen zusätzlich \(\mathrm{Spin}(10)\)-Generatoren & 1648 \\
zusätzlich beide vollständigen Gruppen & 7 \\
volle Gruppen und alle Modenbesetzungen & 1 \\
\end{longtable}
}

Diese bedingten algebraischen Resultate beantworten, welche zusätzlichen
Zugriffe Unterschiede auflösen würden. Sie beweisen keine physische
Verfügbarkeit, Lie-Kontrollierbarkeit oder effiziente Messung. Der
Begriff „volle Algebra`` ersetzt kein Operationsprogramm. {[}S068, §6{]}

